
1.Should transportation funding be based on how much you drive? VMT 


\begin{itemize}
\item FHA , average licensed driver in the US, drives 13-14K miles per year.  


\end{itemize}

\begin{itemize}
\item Oregon running a pilot VMT for some 5000 vehicles and Trucks due to shortage of funding from fuel based tax.  


\end{itemize}

\begin{itemize}
\item  


\end{itemize}

\begin{itemize}
\item Top 3 States (Highest Average Miles Driven) 


\end{itemize}

\begin{itemize}
\item Wyoming — 21,589 miles 


\end{itemize}

\begin{itemize}
\item Indiana — 20,560 miles 


\end{itemize}

\begin{itemize}
\item Mississippi — 19,517 miles 


\end{itemize}

\begin{itemize}
\item Bottom 3 States (Lowest Average Miles Driven) 


\end{itemize}

\begin{itemize}
\item Washington, D.C. — 6,694 miles 


\end{itemize}

\begin{itemize}
\item New York — 9,548 miles 


\end{itemize}

\begin{itemize}
\item Rhode Island — 9,903 miles 


\end{itemize}

Advantage: 


-	It will allow all cars such as EV to be charge  


Geogia been at 17.5K miles / year which means average. Which means you are more likely to driver 2x more as a Georgia than  a NY residence and almost 3X more compare to a Washinton D.C resident 


 


2.Would mileage-based fees disproportionately affect rural or low-income drivers? 


\begin{itemize}
\item I think this will be a burden on low income drivers because a lot of Middle class and upper class people don’t drive that much (Hold meetings online and Remote) 


\end{itemize}

\begin{itemize}
\item 99% of blue collar jobs and workers must commute to site to fix things.  


\end{itemize}

\begin{itemize}
\item  


\end{itemize}

 


3. How should privacy concerns be addressed if mileage tracking is required?  


\begin{itemize}
\item The insurance tracking device  


\end{itemize}

\begin{itemize}
\item  


\end{itemize}


# Mileage-Based Transportation Funding (VMT)
## Verbatim Evidence from the Article

**Source:**
David L. Sjoquist (2023)
*Funding Transportation in Georgia: Vehicle Miles Traveled Tax*
Center for State and Local Finance, Georgia State University

---

## Section 7 — Privacy

### Authority: National Academies (2016), cited by Sjoquist

> “The National Academies (2016) published an analysis of 38 public opinion surveys focused on VMT tax.”

> “However, only four surveys asked questions related to privacy.”

> “Concerns with data privacy was one of the participants’ key objections to VMT tax schemes for which the location and time of travel was collected.”

> “More than half of the respondents indicated that privacy, particularly information on where the respondent drove, was a concern, or they expressed a preference for not collecting such information.”

> “Data collecting that relied solely on an odometer check did not generate undue concern.”

---

### Authority: David L. Sjoquist

> “It is clear that any attempt to implement a VMT tax will require addressing this concern regarding privacy.”

> “One approach to privacy concerns that has been suggested is to have a private firm rather than the government collect the data.”

> “It is worth noting that despite the concerns expressed with VMT tax data, there are widely used tracking systems that collect information on where and when someone drove or have the potential to do so.”

> “For example, E-ZPass systems that are used to charge drivers using toll or express lanes record when a vehicle passes monitoring stations.”

> “General Motors’ OnStar system has access to data on where and when someone drives…”

> “…and insurance company programs such as Allstate DriveWise and Progressive Snapshot not only collect data on driving behavior but also gather information on when and where a vehicle is driven.”

> “It is also possible for entities both private and public to track the location of a person’s cell phone.”

---

## Section 8 — Equity

### Authority: David L. Sjoquist

> “Equity is a subjective concept, and philosophers and other have spilled a lot of ink in a failed attempt to identify a universal operational standard for measuring equity.”

> “Since the policy being considered is a shift from a motor fuel tax to a VMT tax, our concern is with how equity changes as a result of that switch.”

---

### Section 8.1 — Differences by Urban and Rural Residence

> “One can find frequent claims that rural residents are, on average, negatively affected by a shift to a VMT tax. However, that claim does not hold up upon examination.”

> “Miles driven by residents in rural areas are, on average, greater than in urban areas since many residents in rural areas drive long distances to get to work and to access services.”

> “Rural residents also tend to drive larger vehicles that are less fuel efficient.”

> “Thus, rural residents pay relatively more motor fuel taxes than urban residents and are likely to pay more under a VMT tax.”

> “However, the relative difference between what a rural resident would pay and what an urban resident would pay is smaller with a VMT tax.”

---

### Section 8.2 — Differences by Income

> “The VMT tax per mile would be the same regardless of the individual’s income or place of residence.”

> “However, on average the total VMT taxes would be a larger percentage of income for low-income households, which means that VMT taxes are regressive.”

> “However, motor fuel taxes are also regressive and are more regressive than a VMT tax.”

> “Fuel efficiency increases with income since higher income households tend to own newer, more-fuel efficient vehicles.”

> “Thus, switching from a motor fuel tax to a VMT tax is expected to reduce tax regressivity.”

---

### External Authorities Cited in Section 8

> “McMullen, Zhang, and Nakahara (2010) find that rural households benefit more from a switch to a VMT tax as compared to urban households because rural residents own lower-fuel efficiency vehicles.”

> “Schroeckenthaler and Fitzroy (2019) also find that households in urban areas pay slightly more than households in rural areas under a mileage-based road usage charge than under the motor fuel tax.”

> “Weatherford (2011) finds that switching to a VMT tax would shift the tax burden from low to high income households, from older to younger households with children, and from rural to urban households.”

---

## Section 10 — Public Opinion

### Authority: National Academies (2016), cited by Sjoquist

> “There have been several public opinion surveys conducted since 1995 that ask about vehicle miles traveled taxes.”

> “The National Academies (2016) published an extensive meta-analysis of 38 public opinion surveys that, collectively, included 167 unique questions about mileage-based user fees.”

> “The surveys were conducted between 1995 and 2015, all but one of which was conducted in 2005 or later.”

> “Most of the surveys asked fewer than four question about VMT taxes; 21 surveys asked just one question.”

---

### Findings on Support

> “The average percentage of respondents supporting a VMT tax was only 24 percent, with a range across surveys from 8 percent to 50 percent.”

> “There was little difference in support across personal characteristics, other than political affiliation—Democrats were more likely to support a VMT tax than Republicans or independents.”

> “There was a slight upward time trend in the level of support expressed.”

> “Support among participants in two VMT pilot programs was higher, and ranged from 37 percent to 71 percent, with a mean of 51 percent.”

> “This suggests that better understanding of how a VMT tax works might lead to greater support.”

---

### Findings on Fairness Perception

> “It would be difficult to get public support for a VMT tax if the VMT tax was not considered fair.”

> “Only 33 percent and 45 percent of respondents said that a VMT tax is a fair way to raise transportation revenues.”

> “When asked whether a VMT tax was fairer than a motor fuel tax, only 15 percent to 38 percent said yes.”

> “Respondents said that they did not think a VMT tax was fair for rural residents, for those who drive long distances, or for those who have to drive a lot for work.”

> “These opinions are inconsistent with the empirical findings regarding how VMT taxes actually vary across household categories.”

---

## End Note (No Interpretation Added)

All text above:
- Appears verbatim in the article
- Is drawn directly from Sections 7, 8, and 10
- Includes all qualifying language, uncertainty, and caveats
- Adds **no external framing or synthesis**

Judgment is left entirely to the reader.
